% GENERATED FILE. Edit canonical lessons, then run npm run generate:books.
% canonical-source-hash: fnv1a64:a508aee6eaf382da
% canonical-lessons: GE-C521-gedicht, GE-C521-dichter, GE-C521-reim, GE-C521-reimen, GE-C521-schriftsteller, GE-R521-review-514-521

\chapter{Gedichte und Reime}
\label{ch:ge-gedichte-und-reime}
\hlchaptermodality{521}

\begin{chapteropening}
\textbf{By the end of this chapter:} \emph{I can talk about poems and writers.}

Complete the last lesson of chapter 521: I can talk about poems and writers.
\end{chapteropening}

\section[das Gedicht]{das Gedicht — poem}

\label{lesson:GE-C521-gedicht}



\begin{quote}
Before the new one: say the German for non-fiction book, then the German for villain.
\end{quote}

\subsection*{You'll want to know: das Gedicht}
\textbf{das Gedicht} — \textquotedblleft{}poem\textquotedblright{}.

\begin{quote}
Zum Muttertag hat das Kind ein Gedicht gelernt.
\end{quote}

\begin{quote}
\emph{The child learnt a poem for Mother's Day.}
\end{quote}

Plural: \emph{die Gedichte}.

\begin{cousinweb}
\emph{Gedicht} is inherited from Middle High German \emph{getiht}, built on \emph{tihten} \textquotedblleft{}to compose\textquotedblright{} (EtymDB), which is usually traced to Latin \emph{dictāre} \textquotedblleft{}to say often, to dictate\textquotedblright{} (Lewis \& Short).
\end{cousinweb}

\begin{grammarlens}[title={What you've built}]
One more word for telling, explaining and describing at B1.
\end{grammarlens}

\subsection*{Your turn}
\begin{itemize}
\raggedright
  \item \emph{Say it:} \emph{das Gedicht}
  \item \emph{Say it:} \emph{das Gedicht}, once more
  \item \emph{Say it:} a sentence of your own with \emph{das Gedicht}
  \item \emph{Recall:} say the German for non-fiction book, then the German for villain, then say \emph{das Gedicht} again
\end{itemize}

\begin{tcolorbox}[breakable,colback=teal!4,colframe=teal!35!black,title={Before you move on}]
What does das Gedicht mean? (\textquotedblleft{}Poem\textquotedblright{}.) Say it once more.
\end{tcolorbox}
\section[der Dichter]{der Dichter — poet}

\label{lesson:GE-C521-dichter}



\begin{quote}
Before the new one: say the German for villain, then the German for poem.
\end{quote}

\subsection*{You'll want to know: der Dichter}
\textbf{der Dichter} — \textquotedblleft{}poet\textquotedblright{}.

\begin{quote}
Der Dichter las im Café aus seinem neuen Buch vor.
\end{quote}

\begin{quote}
\emph{The poet read from his new book in the café.}
\end{quote}

Plural: \emph{die Dichter}. A woman is \emph{die Dichterin}.

\begin{cousinweb}
\emph{Dichter} is built on \emph{dichten} \textquotedblleft{}to write poetry\textquotedblright{} with \emph{-er}; EtymDB lists \emph{dichten} without a source, and it is usually traced to Latin \emph{dictāre} \textquotedblleft{}to say often, to dictate\textquotedblright{} (Lewis \& Short).
\end{cousinweb}

\begin{grammarlens}[title={What you've built}]
One more word for telling, explaining and describing at B1.
\end{grammarlens}

\subsection*{Your turn}
\begin{itemize}
\raggedright
  \item \emph{Say it:} \emph{der Dichter}
  \item \emph{Say it:} \emph{der Dichter}, once more
  \item \emph{Say it:} a sentence of your own with \emph{der Dichter}
  \item \emph{Recall:} say the German for villain, then the German for poem, then say \emph{der Dichter} again
\end{itemize}

\begin{tcolorbox}[breakable,colback=teal!4,colframe=teal!35!black,title={Before you move on}]
What does der Dichter mean? (\textquotedblleft{}Poet\textquotedblright{}.) Say it once more.
\end{tcolorbox}
\section[der Reim]{der Reim — rhyme}

\label{lesson:GE-C521-reim}



\begin{quote}
Before the new one: say the German for poem, then the German for poet.
\end{quote}

\subsection*{You'll want to know: der Reim}
\textbf{der Reim} — \textquotedblleft{}rhyme\textquotedblright{}.

\begin{quote}
Das Kinderlied hat lustige Reime.
\end{quote}

\begin{quote}
\emph{The children's song has funny rhymes.}
\end{quote}

Plural: \emph{die Reime}.

\begin{cousinweb}
\emph{Reim} is usually traced to Old French \emph{rime} \textquotedblleft{}rhyme\textquotedblright{}, whose own origin is disputed; EtymDB sets the German word beside French \emph{rime}.
\end{cousinweb}

\begin{grammarlens}[title={What you've built}]
One more word for telling, explaining and describing at B1.
\end{grammarlens}

\subsection*{Your turn}
\begin{itemize}
\raggedright
  \item \emph{Say it:} \emph{der Reim}
  \item \emph{Say it:} \emph{der Reim}, once more
  \item \emph{Say it:} a sentence of your own with \emph{der Reim}
  \item \emph{Recall:} say the German for poem, then the German for poet, then say \emph{der Reim} again
\end{itemize}

\begin{tcolorbox}[breakable,colback=teal!4,colframe=teal!35!black,title={Before you move on}]
What does der Reim mean? (\textquotedblleft{}Rhyme\textquotedblright{}.) Say it once more.
\end{tcolorbox}
\section[reimen]{reimen — to rhyme}

\label{lesson:GE-C521-reimen}



\begin{quote}
Before the new one: say the German for poet, then the German for rhyme.
\end{quote}

\subsection*{You'll want to know: reimen}
\textbf{reimen} — \textquotedblleft{}to rhyme\textquotedblright{}.

\begin{quote}
Herz reimt sich auf Schmerz.
\end{quote}

\begin{quote}
\emph{Herz rhymes with Schmerz.}
\end{quote}

Mostly reflexive: \emph{sich reimen auf} + accusative.

\begin{cousinweb}
\emph{reimen} is inherited from Middle High German \emph{rīmen} \textquotedblleft{}to compose, to put into verse, to rhyme\textquotedblright{} (EtymDB).
\end{cousinweb}

\begin{grammarlens}[title={What you've built}]
One more word for telling, explaining and describing at B1.
\end{grammarlens}

\subsection*{Your turn}
\begin{itemize}
\raggedright
  \item \emph{Say it:} \emph{reimen}
  \item \emph{Say it:} \emph{reimen}, once more
  \item \emph{Say it:} a sentence of your own with \emph{reimen}
  \item \emph{Recall:} say the German for poet, then the German for rhyme, then say \emph{reimen} again
\end{itemize}

\begin{tcolorbox}[breakable,colback=teal!4,colframe=teal!35!black,title={Before you move on}]
What does reimen mean? (\textquotedblleft{}To rhyme\textquotedblright{}.) Say it once more.
\end{tcolorbox}
\section[der Schriftsteller]{der Schriftsteller — writer, author}

\label{lesson:GE-C521-schriftsteller}



\begin{quote}
Before the new one: say the German for rhyme, then the German for to rhyme.
\end{quote}

\subsection*{You'll want to know: der Schriftsteller}
\textbf{der Schriftsteller} — \textquotedblleft{}writer, author\textquotedblright{}.

\begin{quote}
Mein Onkel ist Schriftsteller und schreibt Krimis.
\end{quote}

\begin{quote}
\emph{My uncle is a writer and writes crime novels.}
\end{quote}

Plural: \emph{die Schriftsteller}. A woman is \emph{die Schriftstellerin}.

\begin{cousinweb}
\emph{Schriftsteller} is built on \emph{Schrift} \textquotedblleft{}writing\textquotedblright{}, which came through Old High German \emph{scrift} from Latin \emph{scrīptum} (EtymDB), and \emph{Steller}, from \emph{stellen} \textquotedblleft{}to set\textquotedblright{}: one who sets down writing.
\end{cousinweb}

\begin{grammarlens}[title={What you've built}]
One more word for telling, explaining and describing at B1.
\end{grammarlens}

\subsection*{Your turn}
\begin{itemize}
\raggedright
  \item \emph{Say it:} \emph{der Schriftsteller}
  \item \emph{Say it:} \emph{der Schriftsteller}, once more
  \item \emph{Say it:} a sentence of your own with \emph{der Schriftsteller}
  \item \emph{Recall:} say the German for rhyme, then the German for to rhyme, then say \emph{der Schriftsteller} again
\end{itemize}

\begin{tcolorbox}[breakable,colback=teal!4,colframe=teal!35!black,title={Before you move on}]
What does der Schriftsteller mean? (\textquotedblleft{}Writer, author\textquotedblright{}.) Say it once more.
\end{tcolorbox}
\section[(dialogue)]{Words from chapters 514–521 — a review}

\label{lesson:GE-R521-review-514-521}



\begin{quote}
Say the German for melody and rhythm.
\end{quote}

\subsection*{Your turn}
Say these aloud:

\begin{itemize}
\raggedright
  \item melody, rhythm, beat, verse, to hum
  \item choir, orchestra, conductor, to conduct, to compose
  \item stage, to rehearse, stage fright, to perform, appearance
  \item actor, director, audience, to applaud, applause
  \item painting, paintbrush, canvas, easel, watercolour
  \item gallery, to exhibit, sculpture, sculptor, work of art
  \item crime novel, biography, short story, non-fiction book, villain
  \item poem, poet, rhyme, to rhyme, writer
\end{itemize}

\begin{tcolorbox}[breakable,colback=teal!4,colframe=teal!35!black,title={Before you move on}]
Melody? (\textbf{die Melodie}.) Rhythm? (\textbf{der Rhythmus}.) Beat? (\textbf{der Takt}.) Verse? (\textbf{die Strophe}.) To hum? (\textbf{summen}.) Choir? (\textbf{der Chor}.) Orchestra? (\textbf{das Orchester}.) Conductor? (\textbf{der Dirigent}.) To conduct? (\textbf{dirigieren}.) To compose? (\textbf{komponieren}.) Stage? (\textbf{die Bühne}.) To rehearse? (\textbf{proben}.) Stage fright? (\textbf{das Lampenfieber}.) To perform? (\textbf{auftreten}.) Appearance? (\textbf{der Auftritt}.) Actor? (\textbf{der Schauspieler}.) Director? (\textbf{der Regisseur}.) Audience? (\textbf{das Publikum}.) To applaud? (\textbf{klatschen}.) Applause? (\textbf{der Applaus}.) Painting? (\textbf{das Gemälde}.) Paintbrush? (\textbf{der Pinsel}.) Canvas? (\textbf{die Leinwand}.) Easel? (\textbf{die Staffelei}.) Watercolour? (\textbf{das Aquarell}.) Gallery? (\textbf{die Galerie}.) To exhibit? (\textbf{ausstellen}.) Sculpture? (\textbf{die Skulptur}.) Sculptor? (\textbf{der Bildhauer}.) Work of art? (\textbf{das Kunstwerk}.) Crime novel? (\textbf{der Krimi}.) Biography? (\textbf{die Biografie}.) Short story? (\textbf{die Kurzgeschichte}.) Non-fiction book? (\textbf{das Sachbuch}.) Villain? (\textbf{der Bösewicht}.) Poem? (\textbf{das Gedicht}.) Poet? (\textbf{der Dichter}.) Rhyme? (\textbf{der Reim}.) To rhyme? (\textbf{reimen}.) Writer? (\textbf{der Schriftsteller}.)
\end{tcolorbox}
